\begin{lemma}[The H\"older datum at a single ball, from the cited oscillation estimate]
Let $n\ge1$, let $p>1$, let $0<\lambda\le\Lambda$, let $a>0$, let $M\ge0$ and
let $d>0$, and assume the cited H\"older estimate
\noderef{CitedHolderOscillationMinimum}, which is part (a) of (C4), holds in
dimension $n$ at the exponent $p$.  Then there are
$\alpha,R_0,C_1\in\R$ with $0<\alpha\le1$, $R_0>0$ and $C_1\ge0$ --- depending
only on $n,p,\lambda,\Lambda,a,M$ and $d$ --- such that the following holds.

Let $0\le\varepsilon\le1$, let $\Omega\subseteq\R^n$ be open, let
$\Omega'\subseteq\R^n$ satisfy $d\le|x-y|$ for every $x\in\Omega'$ and every
$y\notin\Omega$, and let $x_0\in\Omega'$.  Let $f$ satisfy (F1) of the thesis on
$\Omega$ with the parameters $p,\varepsilon,\lambda,\Lambda,a$
(\noderef{VariationalGrowthF1}), let $u$ be a local minimum for
$F(\cdot;E)=\int_Ef(x,\cdot,\nabla\cdot)\,dx$ in $\Omega$ with weak gradient $g$
(\noderef{IsLocalMinimumOn}), and suppose $|u|\le M$ almost everywhere on
$\Omega$.  Let $0<R\le R_0$ and let $U\colon\R^n\to\R$ be continuous on
$\overline B_R(x_0)$ with $U=u$ at every point of $B_R(x_0)$.  Then
\[
|U(x)-U(x_0)|\le C_1R^{\alpha}\qquad\bigl(x\in\overline B_R(x_0)\bigr).
\]
\end{lemma}
\begin{proof}
\emph{Step 1: the constants.}  Apply \noderef{CitedHolderOscillationMinimum}
with the present $\lambda,\Lambda,a$ and $M$; its hypotheses are $0<\lambda$,
$\lambda\le\Lambda$, $0<a$ and $0\le M$, all present.  It produces
$\alpha,d_\ast$ with $0<\alpha\le1$ and $d_\ast>0$, depending only on
$n,p,\lambda,\Lambda,a,M$, and bound \emph{before} $\varepsilon$ and before all
the data, which is what makes them $\varepsilon$-uniform.  \emph{The subscript
on $d_\ast$ is not decoration}: the cited node names its produced constant
$d$, while the $d$ of the present lemma is the domain distance, and the two are
unrelated --- in particular no inequality between them is available --- so the
cited constant is written $d_\ast$ throughout.  Put
\[
R_0=\min\Bigl\{\frac d2,\,1\Bigr\}>0,
\qquad
C_1=4^{\alpha}\max\Bigl\{2M,\ d_\ast\frac d2\Bigr\}\Bigl(\frac2d\Bigr)^{\alpha} .
\]
Then $R_0>0$ because $d>0$, and $C_1\ge0$ because $4^{\alpha}>0$,
$(2/d)^{\alpha}>0$ and the maximum is at least $2M\ge0$.  Both depend only on
$n,p,\lambda,\Lambda,a,M$ and $d$, as required.

Now fix the data and let $0<R\le R_0$; in particular $R\le d/2$.

\emph{Step 2: the closed ball of radius $d/2$ lies in $\Omega$.}  Let
$y\in\overline B_{d/2}(x_0)$ and suppose $y\notin\Omega$.  Since $x_0\in\Omega'$
the separation hypothesis gives $d\le|x_0-y|$, while $|y-x_0|\le d/2$, so
$d\le d/2$, contradicting $d>0$.  Hence $\overline B_{d/2}(x_0)\subseteq\Omega$,
and a fortiori $B_{d/2}(x_0)\subseteq\Omega$ and, since $R\le d/2$, also
$B_R(x_0)\subseteq\Omega$.

\emph{Step 3: the oscillation over $B_{d/2}$ is at most $2M$.}  The ball
$B_{d/2}(x_0)$ has positive volume (Preamble: the volume of a ball), so the
almost-everywhere filter of the restricted measure is nontrivial and the
essential supremum and infimum of $u$ over it are genuine reals rather than
junk values.  From $|u|\le M$ almost everywhere on $\Omega$ and
$B_{d/2}(x_0)\subseteq\Omega$ we get $u\le M$ and $-M\le u$ almost everywhere
on $B_{d/2}(x_0)$, hence
$\operatorname{ess\,sup}_{B_{d/2}}u\le M$ and
$-M\le\operatorname{ess\,inf}_{B_{d/2}}u$, and therefore
\[
\operatorname*{osc}\bigl(u,B_{d/2}(x_0)\bigr)
=\operatorname{ess\,sup}_{B_{d/2}}u-\operatorname{ess\,inf}_{B_{d/2}}u\le2M
\qquad(\noderef{EssOscOn}).
\]

\emph{Step 4: the cited estimate, scaled.}  Apply the conclusion of
\noderef{CitedHolderOscillationMinimum} at the present $\varepsilon$ (with
$0\le\varepsilon\le1$), the open set $\Omega$, the integrand $f$, the pair
$(u,g)$, the centre $x_0$, the outer radius $d/2$ and the inner radius $R$.  Its
remaining hypotheses are: $\Omega$ open; $f$ satisfying (F1) of the thesis on
$\Omega$ with those parameters (\noderef{VariationalGrowthF1}); $u$ a local
minimum with weak gradient $g$ (\noderef{IsLocalMinimumOn}); $|u|\le M$ almost
everywhere on $\Omega$; $\overline B_{d/2}(x_0)\subseteq\Omega$, by Step 2;
$0<R$; and $R\le d/2$.  It gives
\[
\operatorname*{osc}\bigl(u,B_R(x_0)\bigr)
\le4^{\alpha}\max\Bigl\{\operatorname*{osc}\bigl(u,B_{d/2}(x_0)\bigr),\
d_\ast\frac d2\Bigr\}\Bigl(\frac{R}{d/2}\Bigr)^{\alpha}.
\]
By Step 3 the maximum is at most $\max\{2M,d_\ast d/2\}$, and
$4^{\alpha}>0$ and $(R/(d/2))^{\alpha}\ge0$, so the right-hand side is at most
$4^{\alpha}\max\{2M,d_\ast d/2\}(R/(d/2))^{\alpha}$.  Finally
$R/(d/2)=R\cdot(2/d)$ with both factors nonnegative, so
$(R/(d/2))^{\alpha}=R^{\alpha}(2/d)^{\alpha}$, and therefore
\[
\operatorname*{osc}\bigl(u,B_R(x_0)\bigr)\le C_1R^{\alpha}.
\tag{4}
\]

\emph{Step 5: a continuous function is pointwise trapped by its essential
bounds.}  Write $S=\operatorname{ess\,sup}_{B_R(x_0)}u$ and
$I=\operatorname{ess\,inf}_{B_R(x_0)}u$; as in Step 3 these are genuine reals,
because $B_R(x_0)$ has positive volume and $u$ is essentially bounded there by
$M$.  By definition of the essential supremum and infimum, $I\le u\le S$ almost
everywhere on $B_R(x_0)$, and since $U=u$ at every point of $B_R(x_0)$ the same
holds for $U$: $I\le U\le S$ almost everywhere on $B_R(x_0)$.

We claim $I\le U(x)\le S$ at \emph{every} $x\in B_R(x_0)$.  Suppose
$U(x)>S$ for some $x\in B_R(x_0)$.  The set
$B_R(x_0)\cap U^{-1}\bigl((S,\infty)\bigr)$ is open, being the intersection of
the open ball with the preimage of an open ray under a function continuous on
the ball ($U$ is continuous on $\overline B_R(x_0)$, hence on $B_R(x_0)$), and
it is nonempty, containing $x$; a nonempty open subset of $\R^n$ contains a
ball of positive radius and so has positive volume (Preamble: the volume of a
ball).  But that set is null, because $U\le S$ almost everywhere on
$B_R(x_0)$.  This is a contradiction, so $U\le S$ everywhere on $B_R(x_0)$; the
argument with $(-\infty,I)$ in place of $(S,\infty)$ gives $I\le U$ everywhere
on $B_R(x_0)$.

The bound extends to the closed ball.  The set
$\overline B_R(x_0)\cap U^{-1}\bigl([I,S]\bigr)$ is closed, being the
intersection of the closed ball with the preimage of a closed interval under a
function continuous on the closed ball, and it contains $B_R(x_0)$ by the
previous paragraph.  Hence it contains the closure of $B_R(x_0)$, which is
$\overline B_R(x_0)$ because $R\ne0$.  So $I\le U(x)\le S$ for every
$x\in\overline B_R(x_0)$.

\emph{Step 6: conclusion.}  Let $x\in\overline B_R(x_0)$.  Also
$x_0\in\overline B_R(x_0)$, since $R\ge0$.  By Step 5 both $U(x)$ and $U(x_0)$
lie in $[I,S]$, so
\[
|U(x)-U(x_0)|\le S-I=\operatorname*{osc}\bigl(u,B_R(x_0)\bigr)\le C_1R^{\alpha}
\]
by (4).  That is the assertion.
\end{proof}
